\chapter{17}



Die Unruhe kam mit den Gedanken an Giorgia zurück. Louis` Magen zog sich zusammen. Er legte die Gabel quer über den Teller und verschränkte die Arme. Wieder diese Szene, wie sie vor ihm taumelte, schrie, kämpfte und nach hinten kippte.

Als wäre es gerade eben gewesen, sah er sich am Klavier sitzen. 
Es war Mitte Juli gewesen. Die letzten Abendgäste hatten sich verabschiedet. Die
Überraschung war gelungen. Als seine nächsten Freundinnen und Freunde hatten sie Franco ins Ristorante gelockt. Es sollte eine Geburtstagsüberraschung werden. Kurz nach
Mitternacht. Auch Louis war da. Er hatte sich überwunden. Nach Tagen des
Rückzugs hatte er wieder einmal Zugehörigkeit empfunden, eine leise
Freude sogar. Während er musikalische Wünsche erfüllte, war alles ein
bisschen wie früher. Giorgia direkt anzusehen, verbot er sich. Bis er es nicht mehr schaffte. Sein Blick schweifte zu ihr hinüber und
blieb auf ihr liegen. Er sollte aufstehen und gehen.
Doch dann kam Marco auf ihn zu. Er hielt ihn freundschaftlich an der
Schulter, lachte und blickte ihn auffordernd an.

«Spiel uns noch was, Louis.»

Es klang selbstverständlich, wie oft nach Feierabend, wenn sie noch
zusammensassen. Eine Handvoll Leute, die sich gemütlich eingerichtet
hatte. Zu gemütlich, um nach Hause zu gehen.
Nun klatschte Franco in die Hände.

«Ja, Louis, bitte, bitte.»

Mit einem Mal wusste Louis, der Moment war gekommen, Giorgia den Song zu
spielen. Er hatte ihn über Tage vorbereitet. Auch wenn er nicht überzeugt war, dass sie ihn 
je zu hören bekäme. Damals war er beim ersten Abspielen ausser sich gewesen. Ihm schien,
Alborán habe die Worte von \emph{«Prometo»} bei ihm abgeschrieben.
Treffender ging es nicht.
Er musste versuchen, Giorgia mit diesem Song noch einmal und unmissverständlich alles zu sagen, was ihm wichtig war. Sich ihr in ungeschminkter Verletzlichkeit zeigen. Ihr seine tiefen, wahren Gefühle offenbaren, wie er es bis anhin offensichtlich nicht geschafft hatte. Nicht mit seinen unzähligen Handy-Nachrichten, den
Liebesbezeugungen, bis weit in die Nacht hinein, nicht einmal mit seinem
handgeschriebenen Brief.

Es war gewagt. Er wusste es. Giorgia hatte sich mehrfach und entschieden
für eine endgültige Trennung ausgesprochen.
Ihre Antwort auf seinen Brief und ihre letzte Nachricht hatten
eigentlich keine Fragen mehr offengelassen. Mit steifen Fingern hatte
er das Handy in dieser Nacht gehalten, sich gerade aufgesetzt und
erstmal die Augen geschlossen. Und dann mit einem klitzekleinen
Hoffnungsfunken ihre Worte gelesen.

«Unsere Beziehung, ein Abenteuer, viel schön, viel Drama, mal mehr, mal
weniger, Louis. Alles gesagt. Ich habe nichts mehr, ausser \emph{Birdy’s
Song}, ein musikalischer Abschied für dich, in deiner Sprache, mein
Musiker - er sagt’s, wie ich’s mein. Lass uns Zeit geben, Abstand
schaffen und vielleicht, eines Tages, Freunde werden? Du wirst eine neue
Frau finden, die passt. Frauen mögen dich, Louis. Werde vieles
vermissen, worauf sich deine neue Freundin freuen kann. Ich dank dir
für die Zeit. Bin auch traurig. Die Erinnerungen bleiben. Giorgia»

\qrblock{Birdy and Rhodes \emph{«Let it all go»}}{media/QR-Codes/016_Let_it_all_go_Birdy,_RHODES_Spotify.png}

Sie hatte ihn mit ihren knappen, klaren Worten weggestossen und gedemütigt. Ja, sie hatte \emph{seine} Sprache gewählt. Treffsicher. Und doch, selbst jetzt hielt er sich an ihren letzten drei Worten fest -  \emph{«Bin auch traurig»}. 

\sectionbreak

Louis suchte Giorgias Blick und seine Hände begannen zu zittern. Er fürchtete sich vor dem Versagen seiner Stimme und setzte sich ans Klavier.
Die Freunde rückten ihre Stühle näher. Giorgia sah er links im
Augenwinkel. Seit mehr als einer Woche hatte sie auf stumm geschaltet, 
reagierte auf keine seiner Nachrichten und nahm keine Anrufe entgegen, als
existiere er nicht mehr. Eine Gelegenheit wie diese bot sich vielleicht
nie mehr.
Louis drehte sich zu ihr. Er schaute sie direkt an. Sie hielt den Blick.
«Ich singe für \emph{dich},» sagte er leise. Er wagte kaum zu atmen. Giorgia hob überrascht den Kopf und 
kniff fragend die Augen zusammen. Es lag eine dichte Erwartung in der Luft. Keiner sagte mehr etwas. 
Alle schauten gebannt zwischen Giorgia und Louis hin und her.
Es gab kein Zurück.

\sectionbreak

Giorgia hatte sich kerzengerade aufgesetzt. Unsicher und eine Spur
verlegen rollte sie ihre Hände in ihr weites, schwarzes T-Shirt. Der Raum wurde einzig von einem auf das Klavier gerichteten Spot und 
den letzten brennenden Kerzen beleuchtet, 
die kurz vor dem Erlöschen standen. Es war still.
Der Geruch nach Essen schwebte als Erinnerung an den Abend über ihnen.
Die Klavierklänge begannen sich im Raum zu verteilen. Und Louis sang.

\qrblock{Pablo Alborán \emph{«Prometo»}}{media/QR-Codes/017_Prometo_-_Cata_Acoustic_Sessions_Pablo_Alboràn_Spotify.png}

Sie hatte bis kurz vor dem Ende des Songs durchgehalten. Dann war genug. Noch bevor seine
letzten Silben verklungen waren, stand Giorgia brüsk auf. Der Stuhl
kippte nach hinten. Der Knall traf den Raum hart. Die Hände vor dem
Gesicht rannte sie vornübergebeugt zur Eingangstür, hinaus und ins
Freie. Alle starrten ihr hinterher.

Überwältigung lag drückend schwer im Raum.
Valentina verwarf die Hände und schaute hilfesuchend in die Runde.
